% !PRINTED 226
\weberpage{226}
%<<CONT
ersetzt wird; man hat also nur noch dafür zu sorgen, daß die $\nu$ Werte von (3), die den $\nu$ Reihen entsprechen, voneinander verschieden sind, \gesperrt{um (3) zur Wurzel einer (irreduzibeln) Transformationsgleichung zu machen}\textsuperscript{1)}.

Wir nehmen an, die Funktion $\Phi(x)$ sei entweder eine \gesperrt{gerade} oder eine \gesperrt{ungerade} Funktion von $x$, und machen danach folgende Unterscheidung.

Ist $\Phi(x)$ eine gerade Funktion, so sind unter den Faktoren des Produktes (3) je zwei, nämlich

\weq{4}{\Phi(\mathrm{sn}\,h\,\Omega), \quad \Phi[\mathrm{sn}\,(n-h)\,\Omega]}

einander gleich. Setzen wir also

\weq{5}{\Pi(\Omega) = \prod_{1,\,\frac{n-1}{2}}^{h} \Phi(\mathrm{sn}\,h\,\Omega),}

so ist, wenn $m$ eine beliebige zu $n$ teilerfremde Zahl ist,

\weq{6}{\Pi(m\,\Omega) = \Pi(\Omega)}

[weil unter den $\frac{1}{2}(n-1)$ Zahlen $h\,m$ nicht zwei eine durch $n$ teilbare Summe oder Differenz haben, also, vom Vorzeichen abgesehen, die $\mathrm{sn}\,h\,\Omega$ dieselben sind, wie die $\mathrm{sn}\,h\,m\,\Omega$]. Es ist also $\Pi(\Omega)$ die Wurzel einer Transformationsgleichung. Diese Klasse von Transformationsgleichungen nennen wir \gesperrt{Modulargleichungen}.

Ist $\Phi(x)$ eine ungerade Funktion, so sind die beiden Größen (4) entgegengesetzt und (6) ist nicht mehr allgemein richtig, sondern es ist

\weq{7}{\Pi(m\,\Omega) = \pm\,\Pi(\Omega).}

Daher ist, wenigstens im allgemeinen, nicht mehr $\Pi(\Omega)$, sondern erst $\Pi(\Omega)^2$ Wurzel einer Transformationsgleichung. Diese Art von Transformationsgleichungen nennen wir \gesperrt{Multiplikatorgleichungen}. Um die Fälle kennen zu lernen, in denen $\Pi(\Omega)$ selbst Wurzel einer Transformation ist, muß das Vorzeichen in (7) bestimmt werden.

Dies gelingt auf Grund eines Satzes der Zahlentheorie, der im Bd.\ I, \S\,145 abgeleitet ist.

Der Satz lautet:
%CONT>>

\textsuperscript{1)}~Ausdrücke wie (3) sind auch dann noch Wurzeln von Transformationsgleichungen, wenn $h$ nur zu $n$ teilerfremde Werte annimmt. Solche Transformationsgleichung % [sic] The source prints the plural determiner with a singular noun.
hat man bisher noch wenig benutzt. Wir werden weiterhin ein Beispiel kennen lernen.
